\section{Teorik Gelişim Günlüğü (Theoretical Evolution Log)}

\subsection{Önerme 1.5: Bilişsel Determinizm Sözel Taslağı}
\textbf{Ham Okul Notu (September 2026):} 
Gözlemlenilen kümesine giren her gözlemlenen yani karar verici - ki bunlar her hangi bir unsur olabilir ve sayısı tanımsızdır - kişilikler kümesine gider ve daha sonra karar verme sürecinde beyin bu kümeden veriler alır - daha doğrusu zorunda kalır.
Gk = Gözlemlenilebilirlik kümesi, Gö = Gözlemlenen, Ku = Karar verme süreci fonksiyonu, Ki = Kişilikler kümesi, De = irade.
Ku = (Gö elemanıdır Gk) ise input(Gö elemanıdır Gk) elemanıdır Kj yani, Ku ise değildir De.

\textbf{Sembolik Karşılığı (Formal Notation):}
\[ K_u(p) \implies \Big( (G_{\ddot{o}} \in G_k) \rightarrow \text{input}(G_{\ddot{o}} \in G_k) \in K_i \Big) \implies (K_u \rightarrow \neg D_e) \]

\subsection{Önerme 1.6: Sayısal Motor ve Uzay Bükülmesi (The Curvature Constant)}
\textbf{Gelişim Notu:} Önerme 1.5'i hesaplayan matematiksel bir işlem. Pi sayısı hesaplama hassasiyetinde. Her gözlemlenilen inputla gelen değerler birbirini kombine ederek insan beyni kümesinde noktalar oluşturur. Bazı yüksek şoklu veriler bu alanı daha güçlü büker ($w_n \gg 1$).

\textbf{Matematiksel Karşılığı (Combinatorial Vector Space):}
\[ \Phi: \prod_{n=1}^{N} (w_n \cdot v_n) \longrightarrow P(x_1, x_2, \dots, x_m) \in K_i \]
\[ P_i = \sum_{n=1}^{N} w_n \cdot v_n \]

\subsection{Önerme 1.7: Limit Sonrası ve Sonsuzluk İspatı (Gödelvari Bariyer)}
\textbf{Ham Okul Notu:} (önerme ise zorunlu olası Ls) ise değildir(p ise zorunlu olası Ls) çünkü (olası'dır val(p) elemanıdır önerme kümesi ise 0). Ls burada limit sonrası. Önerme limit sonrası yoktur yani bilinmelidir diyorsa yanlıştır çünkü bir değerin sonsuz olma olasılığı vardır bizde de sonsuzluğu bilememe limiti vardır bu başlı başına limitin kanıtıdır.

\textbf{Kipsel Mantık Karşılığı (Modal Logic Refutation of Rationalism):}
\[ \Big( p \rightarrow \Box \Diamond L_s \Big) \rightarrow \neg \Big( p \rightarrow \Box \Diamond L_s \Big) \quad \because \quad \Big( \Diamond \big( \text{val}(p) \in \mathcal{P}_{set} \big) = 0 \Big) \]
